\section*{Capítulo \ref{ch:g10:functions} --- Funções}

\begin{solution}{exo:g10:functions:1}
$f(0) = 1$; \quad $f(2) = 8 - 6 + 1 = 3$; \quad
$f(-1) = 2 + 3 + 1 = 6$; \quad
$f\!\left(\tfrac12\right) = 2 \times \tfrac14 - \tfrac32 + 1 = \tfrac12 -
\tfrac32 + 1 = 0$.
\end{solution}

\begin{solution}{exo:g10:functions:2}
$f$: sem restrição, \omterm{def:g10:functions:function}{domínio} $\R$.

$g$: exige $x + 4 \neq 0$, \omterm{def:g10:functions:function}{domínio} todos os reais exceto $-4$.

$h$: exige $2x - 6 \geq 0$, isto é, $x \geq 3$: \omterm{def:g10:functions:function}{domínio}
$\intco{3}{+\infty}$.

$k$: $x^2 + 1 \geq 1 > 0$ nunca se anula, \omterm{def:g10:functions:function}{domínio} $\R$.
\end{solution}

\begin{solution}{exo:g10:functions:3}
\emph{\omterm{def:g10:functions:function}{Pré-imagens} de $0$:} $x^2 - 2x = x(x-2) = 0$, logo $0$ e $2$.

\emph{\omterm{def:g10:functions:function}{Pré-imagens} de $3$:} $x^2 - 2x = 3$, isto é, $x^2 - 2x - 3 = 0$,
isto é, $(x-3)(x+1) = 0$ (confira expandindo), logo $3$ e $-1$.

\emph{\omterm{def:g10:functions:function}{Pré-imagens} de $-1$:} $x^2 - 2x + 1 = (x - 1)^2 = 0$, logo a única
\omterm{def:g10:functions:function}{pré-imagem} $1$.
\end{solution}

\begin{solution}{exo:g10:functions:4}
$f(-1) = 1 + 2 + 1 = 4$: sim, $A(-1, 4)$ está no \omterm{def:g10:functions:graph}{gráfico}.
$f(2) = 4 - 4 + 1 = 1$: sim, $B(2, 1)$ também está no \omterm{def:g10:functions:graph}{gráfico}.
\end{solution}

\begin{solution}{exo:g10:functions:5}
\emph{1.} O maior valor do quadro é $4 = g(3)$ (\omterm{def:g10:functions:extrema}{máximo}); o menor é
$-2 = g(0)$ (\omterm{def:g10:functions:extrema}{mínimo}).

\emph{2.} Em $\intcc{-3}{0}$ a \omterm{def:g10:functions:function}{função} decresce, e $-1 < -0.5$, logo
$g(-1) \geq g(-0.5)$.

\emph{3.} Em cada \omterm{def:g10:numbers:interval}{intervalo} de monotonicidade, $g$ assume cada valor no
máximo uma vez, de modo que $g(x) = 0$ tem no máximo uma solução por
\omterm{def:g10:numbers:interval}{intervalo}: no máximo $3$ soluções ao todo. (Aqui $0$ está entre $-2$ e
$1$, entre $-2$ e $4$, e entre $0$ e $4$, de modo que há exatamente
três.)
\end{solution}

\begin{solution}{exo:g10:functions:6}
Sejam $u < v$. Então
\[
f(v) - f(u) = (-2v + 7) - (-2u + 7) = -2(v - u) < 0,
\]
pois $v - u > 0$. Logo $f(v) < f(u)$: $f$ é estritamente \omterm{def:g10:functions:variations}{decrescente} em
$\R$.
\end{solution}

\begin{solution}{exo:g10:functions:7}
Complete o quadrado: $(x+3)^2 = x^2 + 6x + 9$, de modo que
\[
f(x) = x^2 + 6x + 11 = (x + 3)^2 + 2 .
\]
Para todo $x$, $(x+3)^2 \geq 0$, logo $f(x) \geq 2$; e
$f(-3) = 0 + 2 = 2$. O \omterm{def:g10:functions:extrema}{mínimo} de $f$ é $2$, atingido em $x = -3$.
\end{solution}

\begin{solution}{exo:g10:functions:8}
\emph{1.} Comprimento $+$ largura $= 20$, logo o comprimento é $20 - x$ e
\[
A(x) = x(20 - x).
\]
As duas dimensões têm de ser positivas: $0 < x < 20$.

\emph{2.} $A(4) = 64$, $A(8) = 96$, $A(10) = 100$, $A(12) = 96$,
$A(16) = 64$. Os valores sobem até $x = 10$ e depois caem simetricamente:
a área parece ser máxima para $x = 10$, isto é, para um jardim
\emph{quadrado}. (O \cref{ch:g11:quad} demonstra essa conjectura.)
\end{solution}

\begin{solution}{exo:g10:functions:9}
\emph{1.} Para todo $x$: $x^2 + 1 \geq 1 > 0$, logo
$f(x) = \frac{1}{x^2+1}$ é positivo; e, como $x^2 + 1 \geq 1$, inverter
os dois lados (ambos positivos) dá $f(x) \leq 1$.

\emph{2.} \omterm{def:g10:functions:extrema}{Máximo}: $f(0) = 1$ e $f(x) \leq 1$ para todo $x$, logo $1$ é o
\omterm{def:g10:functions:extrema}{máximo}, atingido em $0$. \omterm{def:g10:functions:extrema}{Mínimo}: $f(x) > 0$ sempre, mas nenhum valor de
$x$ atinge $0$ (um quociente de números não nulos é não nulo), e $f$
assume valores tão próximos de $0$ quanto se queira para $x$ grande;
portanto $f$ não tem \omterm{def:g10:functions:extrema}{mínimo}.
\end{solution}

\begin{solution}{exo:g10:functions:10}
\emph{1.} $f(v) - f(u) = v^2 - u^2 = (v-u)(v+u)$. Se $0 \leq u < v$,
então $v - u > 0$ e $v + u > 0$ (soma de um número não negativo com um
positivo), logo $f(v) - f(u) > 0$: $f$ é estritamente \omterm{def:g10:functions:variations}{crescente} em
$\intco{0}{+\infty}$.

\emph{2.} Se $u < v \leq 0$, então $v - u > 0$ ainda, mas agora
$v + u < 0$ (soma de um número não positivo com um negativo), logo
$f(v) - f(u) < 0$: $f$ é estritamente \omterm{def:g10:functions:variations}{decrescente} em
$\intoc{-\infty}{0}$.
\end{solution}

\begin{solution}{pb:g10:functions:1}
\textbf{1.} \omterm{def:g10:functions:function}{Domínio}: $x \neq 0$ (divisão por $x$). $f(1) =
\frac12(1 + 2) = \frac32$; $f(2) = \frac12(2 + 1) = \frac32$;
$f\!\left(\frac32\right) = \frac12\left(\frac32 +
\frac43\right) = \frac{17}{12}$.

\textbf{2.} $\frac12\left(x + \frac2x\right) = \frac32$ dá
$x^2 + 2 = 3x$, isto é, $x^2 - 3x + 2 = (x - 1)(x - 2) = 0$: as
\omterm{def:g10:functions:function}{pré-imagens} de $\frac32$ são $1$ e $2$ --- como a questão 1 já
anunciava.

\textbf{3.} $f(x) = x$ dá $x + \frac2x = 2x$, logo
$\frac2x = x$, $x^2 = 2$: \omterm{pb:g10:functions:1}{pontos fixos} $\sqrt2$ e $-\sqrt2$. O
positivo é a diagonal do quadrado de lado $1$, irracional pelo
problema de fim de semana sobre irracionalidade do volume do
ensino fundamental.

\textbf{4.} $1 \to \frac32 \to \frac{17}{12} \to
\frac{577}{408}$: as \omterm{def:g10:numbers:approx}{aproximações} de $\sqrt2$ de Heron. Sua
receita ``tire a média entre o palpite e $2/$palpite'' é
exatamente \emph{iterar a \omterm{def:g10:functions:function}{função} $f$} --- e o número que o
processo persegue é o \omterm{pb:g10:functions:1}{ponto fixo} de $f$.

\textbf{5.} Os \omterm{pb:g10:functions:1}{pontos fixos} são os cruzamentos do \omterm{def:g10:functions:graph}{gráfico} com a
diagonal $y = x$: realimentar a máquina com suas saídas faz o
ponto caminhar pelo \omterm{def:g10:functions:graph}{gráfico} rumo a esse cruzamento. No problema
de fim de semana sobre os dois termômetros, do volume do ensino
fundamental, a leitura $-40^\circ$ era essa mesma ideia de
cruzar a diagonal, para a \omterm{def:g10:functions:function}{função} de conversão.

\textbf{6.} $f(v) - f(u) = \frac{v - u}{2} + \frac1v - \frac1u
= \frac{v - u}{2} + \frac{u - v}{uv}
= (v - u)\left(\frac12 - \frac{1}{uv}\right)
= \frac{(v - u)(uv - 2)}{2uv}$. Para $0 < u < v \leq \sqrt2$:
$uv < 2$, o parêntese é negativo, $f(v) < f(u)$: \omterm{def:g10:functions:variations}{decrescente}.
Para $\sqrt2 \leq u < v$: $uv > 2$: \omterm{def:g10:functions:variations}{crescente}. Ponto de virada:
$\sqrt2$.

\textbf{7.} \omterm{def:g10:functions:variations}{Decrescente} antes de $\sqrt2$ e \omterm{def:g10:functions:variations}{crescente} depois:
$f$ atinge seu \omterm{def:g10:functions:extrema}{mínimo} em $\sqrt2$, valendo
$f(\sqrt2) = \frac12\left(\sqrt2 + \frac{2}{\sqrt2}\right)
= \sqrt2$. Consequência: toda saída da máquina (com entrada
positiva) é pelo menos $\sqrt2$ --- depois de uma volta da
manivela, os palpites vivem em $\intco{\sqrt2}{+\infty}$.

\textbf{8.} Se $x > \sqrt2$, então $x^2 > 2$, logo
$\frac2x < x$, e a média $f(x)$ entre $x$ e $\frac2x$ é menor
que $x$; e $f(x) \geq \sqrt2$ pela questão 7, com igualdade só
em $x = \sqrt2$. Logo $\sqrt2 < f(x) < x$: cada iteração
melhora estritamente sem ultrapassar --- a sequência da questão
4 desliza até pousar em $\sqrt2$.

\textbf{9.} Em $\intoo{0}{+\infty}$: \omterm{def:g10:functions:variations}{decrescente} desde
$+\infty$ (perto de $0$) até o \omterm{def:g10:functions:extrema}{mínimo} $\sqrt2$ em
$x = \sqrt2$, e depois \omterm{def:g10:functions:variations}{crescente} sem limite.

\textbf{10.} $g(x) = x$ dá $x^2 = x$, $x(x - 1) = 0$: \omterm{pb:g10:functions:1}{pontos
fixos} $0$ e $1$. A partir de $0.9$: $0.81$, $0.656$, $0.430$,
$0.185$ --- deslizando para $0$. A partir de $1.1$: $1.21$,
$1.464$, $2.144$, $4.595$ --- fugindo para o infinito. O \omterm{pb:g10:functions:1}{ponto
fixo} $0$ atrai, o $1$ repele: um fio de diferença na partida,
destinos opostos.

\textbf{11.} $7 \to 22 \to 11 \to 34 \to 17 \to 52 \to 26 \to
13 \to 40 \to 20 \to 10 \to 5 \to 16 \to 8 \to 4 \to 2 \to 1$:
dezesseis passos, com pico em $52$.

\textbf{12.} $27 \to 82 \to 41 \to 124 \to 62 \to 31 \to 94 \to
47 \to 142 \to 71 \to 214 \to \dots$ --- ainda subindo depois de
dez passos, a caminho de um pico de $9\,232$ e de um voo de
$111$ passos. Lição: uma regra de duas linhas pode produzir um
comportamento que ninguém prevê à primeira vista --- números
vizinhos ($26$, $27$, $28$) têm voos radicalmente diferentes.

\textbf{13.} Uma potência de $2$ desce dividindo-se por dois ao
longo de toda a escada: $2^k \to 2^{k-1} \to \dots \to 2 \to
1$. Lá embaixo, $h(1) = 4$, $h(4) = 2$, $h(2) = 1$: o voo entra
no ciclo $4 \to 2 \to 1 \to 4$. Um \omterm{pb:g10:functions:1}{ponto fixo} exigiria
$\frac n2 = n$ (só $n = 0$) ou $3n + 1 = n$ (negativo): nenhum
entre os inteiros positivos --- os voos observados terminam não
em um \omterm{pb:g10:functions:1}{ponto fixo}, mas naquele pequeno \emph{ciclo}.

\textbf{14.} Um contraexemplo seria ou um voo que cresce para
sempre (sem nunca voltar a $1$) ou um segundo ciclo disjunto de
$4 \to 2 \to 1$. A verificação até $10^{20}$ deixa uma
infinidade de números por testar --- exatamente a lição das
regiões do círculo no problema de fim de semana sobre o oráculo
dos palitos, no volume do ensino fundamental: concordar em um
número finito de casos nada prova sobre todos eles.

\textbf{15.} Na máquina de Heron havia \emph{estrutura} --- uma
diferença fatorada (questão 6) revelando a variação, um \omterm{def:g10:functions:extrema}{mínimo},
um cerco --- ao passo que a regra da granizo, que alterna com a
paridade, não oferece alça alguma: ninguém encontrou a
estrutura que a domestique.

\textbf{16.} $\sqrt{x - 3}$: exige $x - 3 \geq 0$:
$\intco{3}{+\infty}$. $\frac{1}{x^2 - 9}$: exige
$x \neq \pm 3$: $\R$ privado de $-3$ e $3$.
$\sqrt{9 - x^2}$: exige $x^2 \leq 9$: $\intcc{-3}{3}$.

\textbf{17.} $f(x) = 3$ dá $x + \frac2x = 6$, logo
$x^2 - 6x + 2 = 0$, isto é, $(x - 3)^2 = 7$:
$x = 3 - \sqrt7$ ou $x = 3 + \sqrt7$ (ambos positivos: duas
\omterm{def:g10:functions:function}{pré-imagens} exatas).

\textbf{18.} Faz sentido físico enquanto a pedra está no ar:
$H(t) = 5t(4 - t) \geq 0$ para $t \in \intcc{0}{4}$.
Completando o quadrado: $H(t) = 20 - 5(t - 2)^2$: altura máxima
$20$ m em $t = 2$ s. Variação: \omterm{def:g10:functions:variations}{crescente} em $\intcc{0}{2}$, de
$0$ a $20$; \omterm{def:g10:functions:variations}{decrescente} em $\intcc{2}{4}$, de volta a $0$.

\textbf{19.} $20t - 5t^2 \geq 15 \iff t^2 - 4t + 3 \leq 0 \iff
(t - 1)(t - 3) \leq 0 \iff t \in \intcc{1}{3}$: dois segundos
inteiros acima de $15$ m.

\textbf{20.} Cada máquina tem um \omterm{def:g10:functions:function}{domínio} (onde a regra é
permitida), uma regra (a fórmula ou o desvio pela paridade) e um
\omterm{def:g10:functions:graph}{gráfico} ou quadro que a exibe por inteiro. A cada uma fizemos as
duas grandes perguntas: \emph{aonde leva a repetição} (\omterm{pb:g10:functions:1}{pontos
fixos}, ciclos, atração --- resolvido para Heron, em aberto para
a granizo) e \emph{como a saída se move com a entrada}
(variação, extremos --- o cume da pedra, o \omterm{def:g10:functions:extrema}{mínimo} de Heron). Os
próximos capítulos aguçam as duas: as funções de referência a
seguir e, mais adiante, as derivadas para medir a variação e os
limites para certificar onde as iterações pousam.
\end{solution}
